\documentclass{article}
\usepackage[T1]{fontenc}
\usepackage{lmodern}
\usepackage{graphicx}
\usepackage{amsmath,amssymb}
\usepackage[a4paper,margin=2cm]{geometry}
\usepackage[absolute,overlay]{textpos}
\setlength{\TPHorizModule}{1mm}
\setlength{\TPVertModule}{1mm}
\usepackage[indonesian]{babel}
\usepackage{hyperref}
\setlength{\emergencystretch}{2em}
% unit: Halaman 7

\begin{document}

% === Halaman 7 (python/native) ===

• Sebuah grup <G, *> dikatakan grup komutatif atau grup abelian (atau 

disingkat abelian saja) jika berlaku aksioma komutatif a * b = b * a untuk 

semua a, b  G.

• <R, +> dan <R, > adalah abelian

• <Z, +> dan <Z, > adalah abelian

• tetapi, <M, >, dengan M adalah himpunan matriks 2 x 2 dengan 

determinan  0 bukan abelian (tanya kenapa?)

Ket: Abelian diambil dari kata “abel”, untuk menghormati Niels Abel, seorang 

Matematikawan Norwegia (1802 – 1829)

Rinaldi Munir/II4021 Kriptografi

7

\end{document}
